\chapter{Estendere il sistema}

\section{Aggiungere una nuova tabella + CRUD}

\begin{enumerate}
  \item Modello in \texttt{webui/backend/models.py}.
  \item Pydantic in \texttt{schemas.py} (\texttt{XxxIn}, \texttt{XxxOut}).
  \item Router in \texttt{routers/xxx.py}. Copia un router esistente
        come template.
  \item Includi il router in \texttt{main.py}:
        \texttt{app.include\_router(xxx.router)}.
  \item Adapter DSL in \texttt{utils/list\_query.py}.
  \item Frontend: nuova route
        \texttt{webui/frontend/src/routes/xxx/+page.svelte},
        riusando \emph{SortableQueryableList}.
  \item Voce in \texttt{+layout.svelte} per la nav.
\end{enumerate}

\section{Migrazione idempotente}

Aggiorna \texttt{db.py::\_apply\_lightweight\_migrations}. Pattern:

\begin{lstlisting}[style=py]
if insp.has_table("my_table") \
        and not has_column("my_table", "new_field"):
    conn.execute(text(
        "ALTER TABLE my_table ADD COLUMN new_field VARCHAR(64)"
    ))
\end{lstlisting}

\section{Aggiungere un nuovo tipo di vincolo}

\paragraph{Per-cella.} Copia il pattern di
\texttt{TeacherUnavailability}: unique (entity\_id, day, hour),
columns state/penalty/reason. Aggiorna \texttt{models.py},
\texttt{schemas.py}, router, \texttt{optimization.py::\_availability\_constraints},
\texttt{routers/monitor.py::\_build\_constraints} (per il tab
Vincoli), \texttt{routers/bulk.py::ENTITY\_UNAV\_MODEL} se vuoi
supportare bulk.

\paragraph{Logico.} Aggiungi \texttt{entity\_type} a
\texttt{LogicalUnavailability}; aggiorna
\texttt{routers/logical.py::ENTITY\_MODEL} e \texttt{ENTITY\_TYPE};
modifica \texttt{routers/monitor.py::\_build\_constraints}; aggiungi
il branch nel solver.

\section{Aggiungere un nuovo step di ottimizzazione}

\begin{lstlisting}[style=py]
def my_new_step(...):
    params = dict(...)
    run_id = create_run("my_step", "Nome leggibile", profile, params)

    def target(rid: int):
        # carica dati, lancia solver, persiste
        update_run(rid, progress=0.5)
        ...
        update_run(rid, solution_id=sid, obj_value=v, metrics=m)
        update_run(rid, progress=1.0)

    start_thread(run_id, target)
    return run_id
\end{lstlisting}

Esponi via \texttt{routers/optimize.py} con un endpoint POST.
Frontend in \texttt{/optimize} per il bottone di lancio + log
streaming via \emph{RunLogPanel} (SSE).

\section{Testing}

Non c'\`e una test suite formale. Per smoke-testare end-to-end:

\begin{itemize}
  \item Backend boot pulito: \texttt{uvicorn backend.main:app}.
  \item Live curl: vedere \S7.2.
  \item Frontend \texttt{vite build} deve passare clean.
  \item Migrazione idempotente: rilancia su DB esistente per
        confermare zero data loss.
\end{itemize}

% =========== Appendice ===========
\appendix

